\chapter{Mở đầu}
\label{chap:chap1-introduce}

\section{Lý do chọn đề tài}

Trong bối cảnh tự động hóa ngày càng phát triển, các hệ thống nhúng đóng vai trò quan trọng trong việc điều khiển và giám sát các quá trình công nghiệp. Đặc biệt, các hệ thống phân loại sản phẩm sử dụng băng truyền được ứng dụng rộng rãi nhằm nâng cao năng suất và độ chính xác trong sản xuất.\\

Để hệ thống băng truyền hoạt động ổn định và đáp ứng tốt yêu cầu phân loại màu sắc, việc thiết kế bộ điều khiển cho động cơ là yếu tố then chốt. Bộ điều khiển PID là phương pháp điều khiển phổ biến, dễ triển khai trên các nền tảng vi điều khiển và phù hợp với các hệ thống nhúng có tài nguyên hạn chế.\\

Xuất phát từ yêu cầu thực tiễn đó, nhóm chúng em lựa chọn đề tài: “Thiết kế bộ điều khiển PID cho hệ thống phân loại màu sắc sử dụng băng truyền”. Đề tài nhằm nghiên cứu, thiết kế và triển khai bộ điều khiển PID trên hệ thống để nâng cao tính ổn định và độ chính xác của quá trình phân loại.

\section{Mục tiêu đề tài}
Mục đích của đề tài là nghiên cứu và thiết kế hệ thống phân loại màu sắc sử dụng băng truyền, trong đó bộ điều khiển PID được áp dụng để điều khiển tốc độ động cơ nhằm đảm bảo hệ thống hoạt động ổn định và chính xác. Thông qua đề tài, nhóm hướng tới việc vận dụng kiến thức đã học về thiết kế hệ thống cũng như điện – điện tử và cảm biến vào mô hình thực tế, đồng thời đánh giá hiệu quả của bộ điều khiển PID trong quá trình phân loại màu.

\section{Đối tượng nghiên cứu}
Đối tượng nghiên cứu của đề tài bao gồm:
\begin{itemize}
    \item Hệ thống phân loại màu sắc sử dụng băng truyền với động cơ DC.
    \item Bộ điều khiển PID giúp điều khiển tốc độ băng truyền.
    \item Cảm biến màu sắc, servo, động cơ DC và quá trình thu nhập, xử lý tín hiệu phục vụ cho quá trình phân loại. 
\end{itemize}

\section{Phạm vi nghiên cứu}
Đề tài được thực hiện trong phạm vi đồ án cuối kỳ của một môn học, tập trung vào việc thiết kế và xây dựng mô hình hệ thống phân loại màu sắc sử dụng băng truyền ở quy mô nhỏ. Nội dung nghiên cứu bao gồm xây dựng cấu trúc hệ thống, thiết kế và hiệu chỉnh bộ điều khiển PID, cũng như đánh giá hoạt động của hệ thống thông qua các thí nghiệm cơ bản. Các nội dung mở rộng và tối ưu hóa nâng cao không thuộc phạm vi của đề tài.

